\documentclass[11pt,a4paper]{article}
\usepackage[utf8]{inputenc}
\usepackage[T1]{fontenc}
\usepackage[portuguese]{babel}
\usepackage{geometry}
\geometry{margin=2.5cm}
\usepackage{listings}
\usepackage{xcolor}
\usepackage{amsmath}
\usepackage{amssymb}
\usepackage{hyperref}
\usepackage{enumitem}
\usepackage{tcolorbox}
\usepackage{tabularx}
\usepackage{booktabs}

\definecolor{codebg}{rgb}{0.95,0.95,0.95}
\definecolor{codegreen}{rgb}{0.1,0.5,0.1}
\definecolor{codepurple}{rgb}{0.5,0.1,0.5}
\definecolor{codegray}{rgb}{0.5,0.5,0.5}

\lstdefinestyle{python}{
    backgroundcolor=\color{codebg},
    commentstyle=\color{codegray}\itshape,
    keywordstyle=\color{codepurple}\bfseries,
    stringstyle=\color{codegreen},
    basicstyle=\ttfamily\small,
    breaklines=true,
    frame=single,
    language=Python,
    showstringspaces=false,
    tabsize=4,
    literate={á}{{\'a}}1 {ã}{{\~a}}1 {é}{{\'e}}1 {í}{{\'i}}1 {ó}{{\'o}}1 {õ}{{\~o}}1 {ú}{{\'u}}1 {ç}{{\c{c}}}1 {Á}{{\'A}}1 {Ã}{{\~A}}1 {É}{{\'E}}1 {Í}{{\'I}}1 {Ó}{{\'O}}1 {Õ}{{\~O}}1 {Ú}{{\'U}}1 {Ç}{{\c{C}}}1 {â}{{\^a}}1 {ê}{{\^e}}1 {ô}{{\^o}}1 {û}{{\^u}}1 {Â}{{\^A}}1 {Ê}{{\^E}}1 {Ô}{{\^O}}1 {Û}{{\^U}}1 {à}{{\`a}}1 {è}{{\`e}}1 {ì}{{\`i}}1 {ò}{{\`o}}1 {ù}{{\`u}}1
}
\lstset{style=python}

\newtcolorbox{objetivos}[1][]{
  colback=green!5!white,
  colframe=green!40!black,
  title=O que você vai aprender nesta página,
  fonttitle=\bfseries,
  #1
}

\newtcolorbox{exercicio}[1]{
  colback=blue!5!white,
  colframe=blue!40!black,
  title=#1,
  fonttitle=\bfseries
}

\newtcolorbox{solucao}{
  colback=yellow!5!white,
  colframe=orange!60!black,
  title=Solução proposta,
  fonttitle=\bfseries
}

\title{\textbf{Data Analysis with Python 2024--2025}\\Parte 6: Machine Learning Types\\\large Caderno de Exercícios}
\author{Tradução e Adaptação: University of Helsinki --- MOOC.fi}
\date{}

\begin{document}

\maketitle

\section*{Nota Importante}
Este material é uma tradução e adaptação baseada no curso \textit{Data Analysis with Python 2024-2025}, oferecido pela \textbf{Universidade de Helsinque (MOOC.fi)}.

\tableofcontents
\newpage

\section{Exercícios - Capítulo 6 (Classificação Naive Bayes)}

\subsection*{Exercício 6.1 -- Classificação de Blobs (blob\_classification)}

\begin{exercicio}{Sumário}
\begin{itemize}
    \item \textbf{Objetivo:} Escrever a função \texttt{blob\_classification(X, y)} que recebe a matriz de \emph{features} \texttt{X} e o vetor de rótulos \texttt{y} e realiza uma classificação do tipo \emph{Naive Bayes Gaussiano}.
    \item \textbf{Procedimento:}
    \begin{enumerate}
        \item Use a função \texttt{train\_test\_split} da biblioteca \texttt{sklearn.model\_selection} para separar os dados em conjunto de treino e conjunto de teste.
        \item Configure \texttt{random\_state=0} para a separação dos dados de forma que o resultado seja sempre determinístico.
        \item Configure \texttt{train\_size=0.75} (ou seja, 75\% dos dados serão usados para treino e 25\% para teste).
        \item Instancie o classificador \texttt{GaussianNB}, faça o fit do modelo com os dados de treinamento e utilize o \texttt{predict} com os dados de teste.
    \end{enumerate}
    \item \textbf{Retorno:} A função deve retornar o \texttt{accuracy\_score} da previsão realizada.
\end{itemize}
\end{exercicio}

\subsection*{Exercício 6.2 -- Classificação de Plantas (plant\_classification)}

\begin{exercicio}{Sumário}
\begin{itemize}
    \item \textbf{Objetivo:} Escrever a função \texttt{plant\_classification()} focada em prever o famoso \textit{Dataset} Iris usando Naive Bayes.
    \item \textbf{Procedimento:}
    \begin{enumerate}
        \item Carregue o conjunto de dados chamando \texttt{sklearn.datasets.load\_iris()}.
        \item Divida os dados chamando \texttt{train\_test\_split} com tamanho de treino equivalente a 80\% (\texttt{train\_size=0.8}) e \texttt{random\_state=0}.
        \item Ajuste (\texttt{fit}) o modelo aos dados utilizando o \texttt{GaussianNB}.
        \item Faça a previsão (\texttt{predict}) sob a parcela de testes.
    \end{enumerate}
    \item \textbf{Retorno:} O escore exato da predição retornado pela função \texttt{accuracy\_score}.
\end{itemize}
\end{exercicio}

\subsection*{Exercício 6.3 -- Classificação de Palavras (word\_classification)}

\begin{exercicio}{Sumário}
\begin{itemize}
    \item \textbf{Objetivo:} Construir um modelo para prever se palavras desconhecidas pertencem ao dicionário finlandês ou ao inglês. Este exercício exige que você processe a limpeza dos dados, conte letras e aplique Validação Cruzada (\textit{Cross-Validation}).
    \item \textbf{Parte 1 - \texttt{get\_features(a)}:} Receba um array unidimensional de palavras e retorne a matriz de \emph{features} no formato \texttt{(n, 29)}. Para cada palavra, a respectiva linha da matriz deve conter o número de ocorrências de cada um dos 29 caracteres possíveis do alfabeto: \texttt{"abcdefghijklmnopqrstuvwxyzäö-"}.
    \item \textbf{Parte 2 - \texttt{contains\_valid\_chars(s)}:} Crie uma função auxiliar que retorne \texttt{True} apenas se a palavra \texttt{s} for composta unicamente pelos 29 caracteres listados acima.
    \item \textbf{Parte 3 - \texttt{get\_features\_and\_labels()}:} Utilize as rotinas que injetam a carga das palavras finlandesas e inglesas (você pode assumir que as listas são carregadas via métodos locais). Transforme o finlandês para letras minúsculas e retire palavras contendo caracteres inválidos (utilizando a Parte 2). Para palavras inglesas, \textbf{remova-as primeiro se começarem com letra maiúscula} (para excluir nomes próprios) e depois aplique as mesmas transformações usadas no finlandês (minúscula e caracteres válidos). Retorne \texttt{(X, y)} sendo finlandês a classe 0 e inglês a classe 1.
    \item \textbf{Parte 4 - \texttt{word\_classification()}:} Extraia \texttt{X, y}. Crie o classificador \texttt{MultinomialNB}. Em seguida, obtenha a matriz de avaliações a partir de \texttt{sklearn.model\_selection.cross\_val\_score} aplicando um \texttt{KFold(n\_splits=5, shuffle=True, random\_state=0)}. Retorne a lista originada pelo \texttt{cross\_val\_score} com as 5 avaliações de acurácia geradas para esse conjunto.
\end{itemize}
\end{exercicio}

\subsection*{Exercício 6.4 -- Detecção de Spam (spam\_detection)}

\begin{exercicio}{Sumário}
\begin{itemize}
    \item \textbf{Objetivo:} Ler e identificar em arquivos brutos compactados se o texto base trata-se de um \emph{Spam} ou de conteúdo limpo (\emph{Ham}).
    \item \textbf{Procedimento:}
    \begin{enumerate}
        \item Abra via módulo nativo do Python \texttt{gzip.open()} os arquivos \texttt{src/ham.txt.gz} e \texttt{src/spam.txt.gz}. Leia o conteúdo das linhas.
        \item Consuma uma quantidade fracionária dos arquivos baseada no parâmetro \texttt{fraction=[0.0, 1.0]}. Você pegará apenas as primeiras linhas correspondentes a esse percentual de arquivo para cada base lida (ex: fraction=0.1 lê os primeiros 10\% de cada arquivo).
        \item Junte as amostras de treino \texttt{ham} seguidas das amostras \texttt{spam} e passe-as em massa pelo \texttt{CountVectorizer().fit\_transform} para converter em contagens vetorizadas de palavras.
        \item Divida os vetores transformados via \texttt{train\_test\_split} usando a amostra de 75\% para treino e \texttt{random\_state=random\_state}. Atribua rótulos $0$ (ham) e $1$ (spam).
        \item Aplique o algoritmo \texttt{MultinomialNB}.
    \end{enumerate}
    \item \textbf{Retorno:} A função \texttt{spam\_detection} deverá retornar uma tupla contendo o \texttt{accuracy\_score} aferido na previsão, a quantidade de elementos testados, e o total absoluto (\texttt{int}) de mensagens testadas que receberam o rótulo incorreto no fim do teste.
\end{itemize}
\end{exercicio}

\end{document}
